\section{Method}
\label{sec:method}

To ask whether \emph{any}
probe-based rule can succeed, we need (i) an evaluation that a rule cannot game
by changing the reasoning it observes, and (ii) a selection procedure that
cannot silently overfit to a lucky operating point. We preregister both.

\subsection{Probe-independent, frozen trajectories}
For each problem we generate exactly \emph{one} main trajectory in a single
request (caps: $16$K tokens for \textsc{MATH}500 and AMC23, $32$K for AIME24)
and freeze it. All probing is then performed \emph{offline} against fixed text
prefixes of that frozen trajectory. Because probes never re-enter the decode,
changing a probe schedule---or a stopping rule---cannot change the underlying
reasoning; every rule is scored against the same trajectories. We build two
offline probe banks: a \emph{dense} bank that re-probes every $64$ tokens with
\simplek{} (a $32$-sample boxed-answer probe), and an \emph{adaptive} bank that
additionally fires event-triggered probes at entropy drops and at
conclusion/reflection/answer markers.

\subsection{Token accounting}
\label{sec:accounting}
For a rule that stops a problem at token position $s$ after consuming probe
output $p$, we charge total decode tokens $T=s+p$ and compare against the frozen
baseline $B$ (the full trajectory, capped at budget). \emph{Net} savings is
$(B-T)/B$ and \emph{gross} savings is $(B-s)/B$; their gap is the probe cost. The
accuracy drop compares the correctness of the committed answer at $s$ against that
of the frozen final answer (it is not an agreement rate between them). Two
quantities therefore move independently, and we track both: \emph{where} the rule
stops (setting the accuracy drop and the main-token saving $s$) and \emph{how much
it probes} to get there (adding $p$). Section~\ref{sec:mechanism} shows that only
the first is intrinsic.

\subsection{The consensus rule space $(W, s)$}
\label{sec:schema}
The consensus signal in the windowed family we study reduces to two
hyperparameters: a \textbf{window size} $W$ (how many of the most recent probes
to inspect) and a \textbf{share threshold} $s$ (what fraction of that window
must agree on a single answer before a stop is allowed). This pair spans the
family to a close approximation: $W{=}1$ is the ``latest-probe'' rule, $s{=}1.0$
is strict unanimity over the last $W$ probes (the CertaIndex-style stop), and an
entropy-over-the-window criterion closely tracks the same agreement and adds no
materially new operating point on our grid. We sweep $W\in\{1,3,5,8,12,16,24,30\}$ and
$s\in\{0.6,0.8,1.0\}$, deliberately including \emph{large} windows so a rule can
demand a long, stable agreement before committing.

Around this core we vary only the operational knobs that govern \emph{when} and
\emph{how} the signal is read: \textbf{probe schedule} (fixed interval
$\in\{64,128,256,512\}$ tokens, or event-triggered probing at
conclusion/entropy/reflection markers); \textbf{maturity} (a minimum-token floor
$\in\{0,512,1024,2048,4096\}$ before any stop is allowed); \textbf{validity} (an
answer-shape filter, rejecting empty/letter answers or only empties); and
\textbf{certainty} (optionally requiring the supporting probes to be flagged
certain). Enumerating this preregistered grid over a fixed-schedule and an
event-triggered family yields $3{,}520$ candidate consensus rules; formal
definitions are in Appendix~\ref{sec:appendix-schema}.

\paragraph{A better signal, swept identically.}
To test whether the failure is early exit \emph{itself} or the consensus
\emph{signal}, we sweep one non-consensus method through the same pipeline,
gates, and token accounting. \textbf{DEER} \citep{deer} stops on the model's own
confidence in a trial answer at a reasoning boundary; its only hyperparameter is
a confidence threshold, which we sweep over $14$ values. We use the stronger
\emph{trial-answer-submit} variant (commit the trial answer directly, no
separate readout) rather than the faithful readout reproduction---our aim here
is not to reproduce prior work but to ask whether \emph{any} rule can clear the
gates. Placing DEER and consensus on one frontier under identical accounting is
what makes the comparison a controlled test of the signal.

\subsection{Models, benchmarks, splits}
Development uses two models---\dsseven{} and \qweneight{}---on three benchmarks
(\textsc{MATH}500, AMC23, AIME24), each at seeds $42/43/44$: $18$
model$\times$benchmark$\times$seed environments. Problems are partitioned
\emph{by problem id} into $60/20/20$ train/dev/test splits
(\textsc{MATH}500 $300/100/100$, AMC23 $24/8/8$, AIME24 $18/6/6$), so no problem
leaks across splits. Confirmation reserves seeds $45/46/47$ plus two unseen
models---\mbox{DeepSeek-R1-Distill-Llama-8B} and \mbox{Distill-Qwen-32B}---all
evaluated on the \emph{test} split only. In one line: \emph{train} generates and
screens candidate rules in-sample, \emph{dev} is the sole selection target the
gates are read on, and \emph{test} (with the two unseen models) is read once for
confirmation after the rules and thresholds are frozen---never for tuning. Where
later sections say ``train+dev,'' they mean a rule was required to pass a gate
in-sample on train and is then measured on dev.

\subsection{Preregistered operating points and gates}
We fix three operating points and their acceptance gates \emph{before} looking
at held-out data (Table~\ref{tab:gates}). Each gate is applied in order:
first the \textbf{total accuracy drop} (macro-averaged over the $18$
environments) must stay under a cap; then the \textbf{total net token saving}
(macro-averaged) must clear a floor; and finally a \textbf{positive-saving
fraction} (psf, the fraction of environments with strictly positive net saving)
floor must be met. A rule is accepted at an operating point only if it satisfies
all three on train+dev. The saving floor is what distinguishes a genuinely
useful early exit from a rule that stays accurate only by almost never stopping.

\begin{table}[t]
  \centering\small
  \begin{tabular}{lccc}
    \toprule
    Operating point & Max total & Min total & psf \\
                    & acc.\ drop & saving    & floor \\
    \midrule
    conservative     & $1.0\pp$ & $10\%$ & $0.80$ \\
    balanced         & $2.0\pp$ & $20\%$ & $0.80$ \\
    token\_efficient & $3.5\pp$ & $30\%$ & $0.70$ \\
    \bottomrule
  \end{tabular}
  \caption{Preregistered acceptance gates. A rule must keep the macro-averaged
  accuracy drop under ``max total acc.\ drop,'' deliver at least ``min total
  saving'' net tokens, and post positive net saving on at least a psf fraction
  of environments. Gates are fixed in advance and \emph{not} relaxed post hoc.}
  \label{tab:gates}
\end{table}

\subsection{Metric resolution}
\label{sec:resolution}
The total accuracy drop is a macro-average over the $18$ environments (each
model$\times$benchmark$\times$seed weighted equally). We span two models, three
benchmarks, and three seeds and weight them equally so that selection rewards a
rule that is robust across benchmarks, models, and randomness, not one that
overfits a single setting. Pooling problems instead would let MATH500 ($100$ dev
problems per seed vs.\ $8$ for AMC23 and $6$ for AIME24) dominate the headline,
so a rule that wins only on MATH500 could pass; equal per-environment weight
forces a selected rule to hold on the harder competition sets too. Macro also
keeps any single problem from swinging the result---one AIME24 problem is
$16.7\pp$ of a $6$-problem split but only ${\sim}1/108$ of the macro drop---while
the three seeds damp the extra per-cell noise of the small sets. The gates are
read on the dev split (the selection target) with train reported alongside;
because \emph{no} consensus rule clears any gate by a wide margin
(Section~\ref{sec:results}), the finding does not rest on a single fragile cell.
Held-out confirmation on the test split and two unseen models provides the
complementary out-of-distribution check.

\subsection{Preregistration commitments}
Three commitments make the result interpretable. \textbf{(C1)} The test split
and both confirmation models are never read during rule selection.
\textbf{(C2)} All three operating points and their gates in
Table~\ref{tab:gates} are fixed before seeing held-out data and are not relaxed
afterwards; if no rule passes a gate, that is reported as the finding rather than
engineered away, and all three points are reported regardless of outcome.
\textbf{(C3)} The frozen rule set is fixed on train+dev and its hash recorded
before any test/confirmation evaluation. Section~\ref{sec:results} reports what
happened when we applied C2.
